\begingroup
\setlength{\abovedisplayskip}{7pt}
\setlength{\belowdisplayskip}{7pt}
\setlength{\abovedisplayshortskip}{3pt}
\setlength{\belowdisplayshortskip}{4pt}

% Printed 152; principal PDF 159; Google PDF 163.
\par\clearpage
\noindent\makebox[0pt][l]{\fontsize{13}{16}\selectfont 152}\hfill{\small\textbf{NACHLASS.}}\hfill\null\par\medskip

\begin{center}
[Berechnung des Vergrößerungsverhältnisses \(n\).]\par\medskip
[6.]
\end{center}
[\(n\) ist das Verhältnis eines Linienelements in der Ebene zum entsprechenden
Element auf dem Ellipsoid.]

Es sei
\[
\begin{aligned}
f'(x)&=a\\
f''(x)&=b\\
f'''(x)&=c\\
f^{\mathrm{IV}}(x)&=d
\end{aligned}
\]
\centerline{usw.}
\medskip
[Man setze
\[
\xi=x-w,
\]
\noindent dann folgt aus
\[
f(x+iy)=f(x-w)+i\lambda:\mathclose{]}
\]
\[
f(x)-\tfrac12byy+\tfrac1{24}dy^4-\tfrac1{720}fy^6\ldots
=f(x)-aw+\tfrac12bww-\tfrac16cw^3\ldots.
\]
[Setzt man]
\[
at=\tfrac12byy-\tfrac1{24}dy^4+\tfrac1{720}fy^6\ldots,
\]
\noindent [so ist]
\[
t=w-\frac b{2a}ww+\frac c{6a}w^3\ldots
\]
\noindent [oder umgekehrt]
\[
w=t+\frac b{2a}tt+\left(\frac{bb}{2aa}-\frac c{6a}\right)t^3\ldots;
\]
\noindent [mithin wird]
\[
w=\frac b{2a}yy+\left(\frac{b^3}{8a^3}-\frac d{24a}\right)y^4
+\left(\frac{b^5}{16a^5}-\frac{b^3c}{48a^4}-\frac{bbd}{48a^3}+\frac f{720a}\right)y^6\ldots
\]
\noindent [und]
\[
\xi=x-\frac b{2a}yy+\frac{aad-3b^3}{24a^3}y^4
-\frac{a^4f-15aabbd-15ab^3c+45b^5}{720a^5}y^6\ldots.
\]
[Da
\[
\begin{aligned}
f'(\xi)&=f'(x)-wf''(x)+\tfrac12wwf'''(x)-\tfrac16w^3f^{\mathrm{IV}}(x)+\ldots\\
&=a-bw+\tfrac12cww-\tfrac16dw^3+\ldots
\end{aligned}
\]

% Printed 153; principal PDF 160; Google PDF 164.
\par\clearpage
\noindent\null\hfill{\small\textbf{KONFORME ABBILDUNG DES SPHÄROIDS IN DER EBENE.}}\hfill\makebox[0pt][r]{\fontsize{13}{16}\selectfont 153}\par\medskip

\noindent ist, so wird]
\[
\begin{aligned}
f'(\xi)=a\biggl\{1-\tfrac12\frac{bb}{aa}yy
 &+\left(\tfrac1{24}\frac{bd}{aa}+\tfrac18\frac{bbc}{a^3}-\tfrac18\frac{b^4}{a^4}\right)y^4\\
 &-\left(\tfrac1{720}\frac{bf}{aa}+\tfrac1{48}\frac{bcd}{a^3}
 -\tfrac1{12}\frac{b^4c}{a^5}+\tfrac1{16}\frac{b^6}{a^6}\right)y^6\ldots\biggr\}.
\end{aligned}
\]
[Nach S.\ 144 ist
\[
\begin{aligned}
f(\xi)&=f(x)-\tfrac12byy+\tfrac1{24}dy^4\ldots\\
\lambda&=ay-\tfrac16cy^3+\tfrac1{120}ey^5\ldots;
\end{aligned}
\]
\noindent setzt man
\[
\begin{aligned}
\frac{\partial f(\xi)}{\partial x}={]}&\ a\left(1-\tfrac12\frac ca yy+\tfrac1{24}\frac ea y^4\ldots\right)=u\\
\left[\frac{\partial\lambda}{\partial x}={}\right]&\ a\left(\frac ba y-\tfrac16\frac da y^3+\tfrac1{120}\frac fa y^5\ldots\right)=v,
\end{aligned}
\]
\noindent [so ist]
\[
\begin{aligned}
\frac{u}{f'(\xi)}&=\frac{\cos c}{n}=1+\left(-\tfrac12\frac ca+\tfrac12\frac{bb}{aa}\right)yy\\
 &\quad+\left(\tfrac1{24}\frac ea-\tfrac1{24}\frac{bd}{aa}-\tfrac38\frac{bbc}{a^3}
 +\tfrac38\frac{b^4}{a^4}\right)y^4\ldots\\
\frac{v}{f'(\xi)}&=\frac{\sin c}{n}=\frac ba y
-\left(\tfrac16\frac da-\tfrac12\frac{b^3}{a^3}\right)y^3\ldots,
\end{aligned}
\]
\noindent [wobei \(c\) wieder die Meridiankonvergenz bezeichnet. Daraus ergibt sich]
\[
\frac1{nn}=1+\left(-\frac ca+2\frac{bb}{aa}\right)yy
+\left(\tfrac1{12}\frac ea-\tfrac5{12}\frac{bd}{aa}+\tfrac14\frac{cc}{aa}
-\tfrac54\frac{bbc}{a^3}+2\frac{b^4}{a^4}\right)y^4\ldots
\]
\noindent [und]
\[
2\log n=\left(\frac ca-2\frac{bb}{aa}\right)yy
+\left(-\tfrac1{12}\frac ea+\tfrac5{12}\frac{bd}{aa}+\tfrac14\frac{cc}{aa}
-\tfrac34\frac{bbc}{a^3}\right)y^4\ldots.
\]
\noindent [Der Logarithmus ist hier wie in den folgenden Artikeln der hyperbolische Logarithmus.]

\begin{center}[7.]\end{center}
[Es ist auch]
\[
\log n-ic=\log f'(\xi)-\log f'(x+iy),
\]
\noindent [mithin]
\[
\log n=\log f'(\xi)-\text{Realteil }\log f'(x+iy).
\]
[Nun ist]
\[
\log f'(\xi)=\log a-\frac{bb}{2aa}yy+\frac{aabd+3abbc-6b^4}{24a^4}y^4\ldots,
\]
\[
[f'(x+iy)=a\left(1+i\frac ba y-\tfrac12\frac ca yy-\tfrac16i\frac da y^3
+\tfrac1{24}\frac ea y^4\ldots\right)
\]
\noindent\textsc{ix.}\hfill 20\hspace*{.12\linewidth}\par

% Printed 154; principal PDF 161; Google PDF 165.
\par\clearpage
\noindent\makebox[0pt][l]{\fontsize{13}{16}\selectfont 154}\hfill{\small\textbf{NACHLASS.}}\hfill\null\par\medskip

\[
\begin{aligned}
\log f'(x+iy)=\log a
 &+\tfrac12\left(\frac{bb}{aa}-\frac ca\right)yy
 +\left(\frac{ae-4bd+3cc}{24aa}-\frac{(bb-ac)^2}{4a^4}\right)y^4\ldots\\
 &+i\left(\frac ba y-\left(\tfrac16\frac da-\tfrac12\frac{bc}{aa}
 +\tfrac13\frac{b^3}{a^3}\right)y^3\ldots\right)\mathclose{]}
\end{aligned}
\]
\[
\text{Realteil }\log f'(x+iy)=\log a+\frac{bb-ac}{2aa}yy
+\frac{a^3e-4aabd-3aacc+12abbc-6b^4}{24a^4}y^4\ldots.
\]
\noindent Also
\[
\begin{aligned}
\log n&=\frac{ac-2bb}{2aa}yy-\frac{aae-5abd-3acc+9bbc}{24a^3}y^4\ldots
 \quad\text{[vgl. S.\ 153;}\\
c&=\frac ba y-\frac{aad-3abc+2b^3}{6a^3}y^3\ldots,
 \quad\text{vgl. S.\ 147.}
\end{aligned}
\]
\noindent Setzt man in die erste Gleichung für \(a=f'(x)\), \(b=f''(x)\) usw. die
Werte von S.\ 144 ein, so erhält man]
\[
\begin{aligned}
\log n={}&\frac{(1-ee\sin\varphi^2)^2}{2aa(1-ee)}yy\\
 &-\frac{(1-ee\sin\varphi^2)^3}{12a^4(1-ee)^3}
 \{1-3ee+(ee+15e^4)\sin\varphi^2-14e^4\sin\varphi^4\}y^4\ldots
\end{aligned}
\]
\noindent [wobei \(a\) wieder die große Halbachse und \(e\) die Exzentrizität bezeichnet].

\begin{center}[8.]\end{center}
Oder setzt man
\[
\begin{aligned}
\frac1{f'(x)}&=\theta(x)=\alpha=\frac1a\\
\theta'(x)&=\beta=-\frac b{aa}=-\sin\varphi\\
\theta''(x)&=\gamma=\frac{2bb-ac}{a^3}\\
\theta'''(x)&=\delta=\frac{-6b^3+6abc-aad}{a^4}\\
\theta^{\mathrm{IV}}(x)&=\varepsilon
=\frac{24b^4-36abbc+6aacc+8aabd-a^3e}{a^5}\qquad\text{usw.,}
\end{aligned}
\]
\noindent [so wird]
\[
\begin{aligned}
\log n&=-\frac{\gamma}{2\alpha}yy+
 \frac{\alpha\alpha\varepsilon-3\alpha\beta\delta-3\alpha\gamma\gamma+3\beta\beta\gamma}{24\alpha^3}
 y^4\ldots,\\
n&=1-\frac{\gamma}{2\alpha}yy+
 \frac{\alpha\alpha\varepsilon-3\alpha\beta\delta+3\beta\beta\gamma}{24\alpha^3}
 y^4\ldots.
\end{aligned}
\]

\begin{center}[9.]\end{center}
Man hat auch [wenn man \(\dfrac{1}{f'(\xi)}=\theta(\xi)\) setzt]
\[
\frac{\partial\partial\log n}{\partial x^2}+\frac{\partial\partial\log n}{\partial y^2}
=-\frac{\theta''(\xi)}{nn\theta(\xi)}.
\]
Ist nun
\[
\frac{\theta''(\xi)}{\theta(\xi)}=-A-Byy-Cy^4-\ldots,
\]

% Printed 155; principal PDF 162; Google PDF 166.
\par\clearpage
\noindent\null\hfill{\small\textbf{KONFORME ABBILDUNG DES SPHÄROIDS IN DER EBENE.}}\hfill\makebox[0pt][r]{\fontsize{13}{16}\selectfont 155}\par\medskip

\noindent so ist \(\left[\dfrac{\partial A}{\partial x}=A',
\dfrac{\partial A'}{\partial x}=A'',\ \text{usw. gesetzt}\right]\)
\[
\begin{aligned}
\log n={}&\tfrac12Ayy+\tfrac1{12}(B-AA-\tfrac12 A'')y^4\\
 &+\bigl(\tfrac1{30}C-\tfrac1{360}B''-\tfrac7{180}AB+\tfrac1{45}A^3
 +\tfrac1{180}A'A'+\tfrac1{120}AA''+\tfrac1{720}A^{\mathrm{IV}}\bigr)y^6\ldots.
\end{aligned}
\]

\begin{center}[10.]\end{center}
[Es ist]
\[
\begin{gathered}
\log n=N\\
\frac{\partial\partial N}{\partial x^2}+\frac{\partial\partial N}{\partial y^2}
=\frac{(1-ee\sin\Phi^2)^2}{aann(1-ee)}
\end{gathered}
\]
\noindent oder
\[
\frac{n\partial\partial n}{\partial x^2}+\frac{n\partial\partial n}{\partial y^2}
-\left(\frac{\partial n}{\partial x}\right)^2-\left(\frac{\partial n}{\partial y}\right)^2
=\frac{(1-ee\sin\Phi^2)^2}{aa(1-ee)}.
\]

\begin{center}
[Beziehungen zwischen \(x,y\) und \(\xi,\lambda\).]\par\medskip
[11.]
\end{center}
[Die Umkehrung der Reihen für \(\xi\) und \(\lambda\) in Art.\ 6 ergibt Folgendes, da nach Art.\ 8
\[
a=\frac1\alpha,\quad b=-\frac{\beta}{\alpha\alpha},\quad
c=-\frac{\gamma}{\alpha\alpha}+\frac{2\beta\beta}{\alpha^3},\quad
d=-\frac{\delta}{\alpha\alpha}+\frac{6\beta\gamma}{\alpha^3}
-\frac{6\beta^3}{\alpha^4},\quad\text{usw. ist,]}
\]
\[
\begin{aligned}
x&=\xi-\tfrac12\alpha\beta\lambda\lambda
+\tfrac1{24}(\alpha^3\delta-2\alpha\alpha\beta\gamma-5\alpha\beta^3)\lambda^4\ldots\\
y&=\alpha\lambda-\tfrac16(\alpha\alpha\gamma-2\alpha\beta\beta)\lambda^3\ldots.
\end{aligned}
\]

\begin{center}[12.]\end{center}
Setzt man
\[
\int\mathrm e^{-\vartheta(t)}\,\mathrm dt=f(t),\qquad
\text{[also nach Art.\ 8: }\theta(t)=\mathrm e^{\vartheta(t)}\text{]}
\]
\noindent und
\[
\begin{array}{l@{\quad}l@{\quad}l}
\vartheta'(x)=\alpha_1&\text{[so gilt]:}&f''(x)=(-\alpha_1)f'(x)\\
\vartheta''(x)=\beta_1&&f'''(x)=(\alpha_1\alpha_1-\beta_1)f'(x)\\
\vartheta'''(x)=\gamma_1&&
 f^{\mathrm{IV}}(x)=(-\alpha_1^3+3\alpha_1\beta_1-\gamma_1)f'(x)\\
\vartheta^{\mathrm{IV}}(x)=\delta_1&&
 f^{\mathrm V}(x)=(\alpha_1^4-6\alpha_1\alpha_1\beta_1+4\alpha_1\gamma_1
 +3\beta_1\beta_1-\delta_1)f'(x)\\
\mathopen{[}\vartheta^{\mathrm V}(x)=\varepsilon_1&&
 f^{\mathrm{VI}}(x)=(-\alpha_1^5+10\alpha_1^3\beta_1-10\alpha_1\alpha_1\gamma_1
 -15\alpha_1\beta_1\beta_1+5\alpha_1\delta_1\\
&&\multicolumn{1}{r}{{}+10\beta_1\gamma_1-\varepsilon_1)f'(x)]}\\
\qquad\text{usw.}&&\multicolumn{1}{c}{\text{usw.,}}
\end{array}
\]
\noindent\hfill \(20^{*}\)\hspace*{.17\linewidth}\par
% Printed 156; principal PDF 163; Google PDF 167.
\par\clearpage
\noindent\makebox[0pt][l]{\fontsize{13}{16}\selectfont 156}\hfill{\small\textbf{NACHLASS.}}\hfill\null\par\medskip

\noindent [und die Reihen für \(\xi\) und \(\lambda\) in Art.\ 6 gehen über in]
\[
\begin{aligned}
\xi&=x+\tfrac12\alpha_1yy+
  (\tfrac1{12}\alpha_1^3+\tfrac18\alpha_1\beta_1-\tfrac1{24}\gamma_1)y^4\\
 &\quad+(\tfrac1{45}\alpha_1^5+\tfrac5{72}\alpha_1^3\beta_1
 -\tfrac1{144}\alpha_1\alpha_1\gamma_1+\tfrac1{48}\alpha_1\beta_1\beta_1
 -\tfrac1{144}\alpha_1\delta_1-\tfrac1{72}\beta_1\gamma_1
 +\tfrac1{720}\varepsilon_1)y^6\ldots\\
\lambda&=yf'(x)\{1-\tfrac16(\alpha_1\alpha_1-\beta_1)yy+
  \tfrac1{120}(\alpha_1^4-6\alpha_1\alpha_1\beta_1+4\alpha_1\gamma_1
  +3\beta_1\beta_1-\delta_1)y^4\ldots\}.
\end{aligned}
\]

\begin{center}
[13.]\par\medskip
[Berechnung der ebenen rechtwinkligen Koordinaten aus der geografischen Breite und Länge.]
\end{center}
Bei der umgekehrten Aufgabe [die Koordinaten \(x,y\) aus der Breite \(\Phi\) und
der Länge \(\lambda\) zu berechnen] bezeichnen wir mit \(\mathrm g\) die Umkehrfunktion
von \(f\), sodass \(\mathrm g(f(x))=x\) ist, und setzen \(\mathrm F(\Phi)=\mathrm F\).

Man hat dann
\[
\begin{aligned}
\mathopen{[}x+iy&=\mathrm g(\mathrm F+i\lambda)\\
x&=\mathrm g(\mathrm F)-\frac{\lambda\lambda}{2}\mathrm g''(\mathrm F)
 +\frac{\lambda^4}{24}\mathrm g^{\mathrm{IV}}(\mathrm F)\ldots\\
y&=\lambda\mathrm g'(\mathrm F)-\frac{\lambda^3}{6}\mathrm g'''(\mathrm F)
 +\frac{\lambda^5}{120}\mathrm g^{\mathrm V}(\mathrm F)\ldots]\\
\mathrm g(\mathrm F)&=\mathrm X,
\end{aligned}
\]
\noindent [wenn \(\mathrm X\) nun die Abszisse bezeichnet, die zum Schnittpunkt des
Breitenkreises \(\Phi\) mit dem Hauptmeridian gehört. Nach Art.\ 1 ist daher
\(f(\mathrm X)=\mathrm F(\Phi)\), also
\(\dfrac{\mathrm d\mathrm X}{\mathrm d\mathrm F}=\dfrac1{f'(\mathrm X)}\). Damit ergibt sich]
\[
\begin{gathered}
\mathrm g'(\mathrm F)=\frac1{f'(\mathrm X)}=\frac{ac}{p}=h,\\
c=\cos\Phi,\qquad s=\sin\Phi,\qquad p=\surd(1-ee\sin\Phi^2)
\end{gathered}
\]
[\(a\) große Halbachse und \(e\) Exzentrizität].

Für die weitere Differentiation [gilt nach S.\ 144]
\[
\frac{\mathrm d\Phi}{\mathrm d\mathrm F}=\frac{ppc}{1-ee}
 \left[=c+\delta c^3;\qquad\frac{ee}{1-ee}=\delta\right]
\]
\noindent oder, wenn man die höheren Potenzen von \(ee\) vernachlässigt,
\(\dfrac{\mathrm d\Phi}{\mathrm d\mathrm F}=c+eec^3\). Mithin

% Printed 157; principal PDF 164; Google PDF 168.
\par\clearpage
\noindent\null\hfill{\small\textbf{KONFORME ABBILDUNG DES SPHÄROIDS IN DER EBENE.}}\hfill\makebox[0pt][r]{\fontsize{13}{16}\selectfont 157}\par\medskip

\[
\begin{aligned}
\frac{\mathrm ds}{\mathrm d\mathrm F}
 &=\frac{pp(1-ss)}{1-ee}\left[=cc+\delta c^4,\qquad
 \frac{\mathrm dc}{\mathrm d\mathrm F}=-\frac{ppsc}{1-ee}=-s(c+\delta c^3)\right];\\
\frac{\mathrm dh}{\mathrm d\mathrm F}&=-sh.
\end{aligned}
\]
Gilt also \(\mathrm g^n(\mathrm F)=hu\), wobei \(u\) eine rationale Funktion von \(c\) ist, so ergibt sich
\[
\left[\mathrm g^{n+1}(\mathrm F)=-sh\left\{u+\frac{\mathrm du}{\mathrm dc}(c+\delta c^3)\right\}\right.
\]
\noindent und, wenn man die höheren Potenzen von \(ee\) vernachlässigt,]
\[
\mathrm g^{n+1}(\mathrm F)=-sh\left\{u+\frac{\mathrm du}{\mathrm dc}(c+eec^3)\right\};
\]
\noindent ist aber
\[
\mathrm g^n\mathrm F=shu,
\]
\noindent wobei \(u\) eine rationale Funktion von \(c\) ist, so ergibt sich
\[
\left[\mathrm g^{n+1}(\mathrm F)=-h\left\{u(1-2cc-\delta c^4)
 +\frac{\mathrm du}{\mathrm dc}(1-cc)(c+\delta c^3)\right\}\right.
\]
\noindent und, wenn man die höheren Potenzen von \(ee\) vernachlässigt,]
\[
\mathrm g^{n+1}(\mathrm F)=-h\left\{u(1-2cc-eec^4)
 +\frac{\mathrm du}{\mathrm dc}(1-cc)(c+eec^3)\right\}.
\]
Bezeichnet man [also] der Kürze halber
\(\dfrac{ee}{1-ee}=\dfrac{2\omega-1}{(\omega-1)^2}\) mit \(\delta\),
[wobei \(\omega\) die Abplattung ist,] so ergeben sich die aufeinanderfolgenden Ableitungen
\[
\begin{aligned}
\mathopen{[}\mathrm g'(\mathrm F)&=\phantom{+}h\\
\mathrm g''(\mathrm F)&=-hs\\
\mathrm g'''(\mathrm F)&=+h(1-2cc-\delta c^4)\\
\mathrm g^{\mathrm{IV}}(\mathrm F)&=-hs(1-6cc-9\delta c^4-4\delta\delta c^6)\\
\mathrm g^{\mathrm V}(\mathrm F)&={]}+h(1-20cc+(24-58\delta)c^4+(72\delta-64\delta\delta)c^6
 +(77\delta\delta-24\delta^3)c^8+28\delta^3c^{10})
\end{aligned}
\]
\centerline{usw.}
\medskip

Folglich [wenn nun sowohl \(x\) als auch \(\mathrm X\) von einem bestimmten Ausgangspunkt
an nach Süden positiv gezählt werden und \(y\) dasselbe Vorzeichen wie \(\lambda\) hat]
\[
\begin{aligned}
x&=\mathrm X-\frac{asc}{2\surd(1-eess)}\lambda\lambda
 +\frac{asc}{24\surd(1-eess)}(1-6cc-9\delta c^4-4\delta\delta c^6)\lambda^4\ldots\\
y&=\frac{ac}{\surd(1-eess)}\lambda
 -\frac{ac}{6\surd(1-eess)}(1-2cc-\delta c^4)\lambda^3\\
 &\quad+\frac{ac}{120\surd(1-eess)}(1-20cc+(24-58\delta)c^4\ldots)\lambda^5\ldots
\end{aligned}
\]

% Printed 158; principal PDF 165; Google PDF 169.
\par\clearpage
\noindent\makebox[0pt][l]{\fontsize{13}{16}\selectfont 158}\hfill{\small\textbf{NACHLASS.}}\hfill\null\par\medskip

\noindent [oder auch]
\[
\begin{aligned}
x&=\mathrm X-\frac{asc}{2p}\lambda\lambda-\frac{asc}{24p(1-ee)^2}
 (5-ee-(6+6ee)ss+(9ee+3e^4)s^4-4e^4s^6)\lambda^4\ldots\\
y&=\frac{ac}{p}\lambda+\frac{ac}{6p(1-ee)}(1-2ss+ees^4)\lambda^3\ldots
\end{aligned}
\]
\noindent [oder]
\[
y=\frac{ac}{p}\sin\lambda+\frac{ac^3}{3p(1-ee)}
 (1-\tfrac12ee-\tfrac12eess)\sin\lambda^3\ldots,
\]
\noindent wofür man auch schreiben kann [wenn die höheren Potenzen von \(ee\) vernachlässigt werden]
\[
y=\frac{ac}{p}\sin\lambda+\frac{ac^3}{3\surd(1-ee)}\sin\lambda^3.
\]
[Vernachlässigt man dieselben Terme, so ergibt sich]
\[
\begin{aligned}
x&=\mathrm X-\frac{2asc}{p}\sin\tfrac12\lambda^2
 \left(1+cc(2+3eecc)\sin\tfrac12\lambda^2\right)\\
 &=\mathrm X-\frac{2asc}{p}\sin\tfrac12\lambda^2
 \left(1+2cc\frac{p^3}{(1-ee)^{\frac32}}\sin\tfrac12\lambda^2\right).
\end{aligned}
\]

\begin{center}
[14.]\par\medskip
Berechnung der Meridiankonvergenz aus den geografischen Koordinaten.
\end{center}
[Die Meridiankonvergenz \(c\) ist auch gleich dem Winkel, den im Punkt \(\Phi,\lambda\)
der Meridian mit der Kurve bildet, deren Abbildung in der Ebene eine Parallele
zur \(x\)-Achse ist. Dann gilt
\[
\begin{aligned}
\operatorname{tang}c&=\frac{-\lambda\mathrm g''(\mathrm F)+\tfrac16\lambda^3\mathrm g^{\mathrm{IV}}(\mathrm F)\ldots}
 {\mathrm g'(\mathrm F)-\tfrac12\lambda\lambda\mathrm g'''(\mathrm F)\ldots}\\
 &=-\lambda\frac{\mathrm g''(\mathrm F)}{\mathrm g'(\mathrm F)}
 +\tfrac16\lambda^3\left(\frac{\mathrm g^{\mathrm{IV}}(\mathrm F)}{\mathrm g'(\mathrm F)}
 -3\frac{\mathrm g''(\mathrm F)\mathrm g'''(\mathrm F)}{(\mathrm g'(\mathrm F))^2}\right)\ldots\mathclose{]}\\
\operatorname{tang}c&=\lambda\sin\Phi+\tfrac13\lambda^3\sin\Phi
 (1+3\delta\cos\Phi^4+2\delta\delta\cos\Phi^6)\ldots
\end{aligned}
\]
\noindent [oder]
\[
\begin{aligned}
\operatorname{tang}c&=\operatorname{tang}\lambda\sin\Phi
 \left\{1+\tfrac13\operatorname{tang}\lambda^2\cos\Phi^4\,.\,ee
 \frac{3-ee-2ee\sin\Phi^2}{(1-ee)^2}\ldots\right\}\\
 &=\operatorname{tang}\lambda\sin\Phi
 \left\{1+\tfrac13\operatorname{tang}\lambda^2\cos\Phi^4\,.\,ee
 \frac{3(1-ee)+2ee\cos\Phi^2}{(1-ee)^2}\ldots\right\}
\end{aligned}
\]
\noindent [oder auch]
\[
c=\lambda\sin\Phi+\tfrac13\lambda^3\sin\Phi\cos\Phi^2
 (1+3\delta\cos\Phi^2+2\delta\delta\cos\Phi^4)\ldots.
\]

% Printed 159; principal PDF 166; Google PDF 170.
\par\clearpage
\noindent\null\hfill{\small\textbf{KONFORME ABBILDUNG DES SPHÄROIDS IN DER EBENE.}}\hfill\makebox[0pt][r]{\fontsize{13}{16}\selectfont 159}\par\medskip

\begin{center}
[15.]\par\medskip
[Reduktion des sphäroidischen Azimuts auf das ebene Azimut.]
\end{center}
Die Korrektur der beobachteten Azimute ist
\[
-\alpha(x'-x)\eta+\beta(x'-x)\eta^3+\gamma(y'-y)\eta\eta,
\]
\noindent wo \(\eta=\tfrac13(y'+2y)\).

[\(x,y\) sind die ebenen Koordinaten des Beobachtungspunkts, \(x'y'\) die
Koordinaten des Punkts, zu dem das Azimut bestimmt wird. Die \(x\)-Achse ist
dabei nach Süden, die \(y\)-Achse nach Westen positiv.]
\[
\begin{aligned}
\alpha&=\frac1{2aa(1-ee)qq}\\
\beta&=\frac{1-3ee+(ee+15e^4)\sin\varphi^2-14e^4\sin\varphi^4}{6a^4(1-ee)^3q^3}\\
\gamma&=\frac{ee\sin2\Phi}{2a^3(1-ee)^2q^{\frac52}}\\
\mathopen{[}q&=\frac1{1-ee\sin\varphi^2},
\end{aligned}
\]
\noindent oder in Bogensekunden und für Briggs'sche Logarithmen]
\[
\alpha=\frac{\mathfrak A}{qq},\qquad
\beta=\frac{\mathfrak B}{q^3}\{1+\mathfrak C\sin\varphi^2-\mathfrak D\sin\varphi^4\},\qquad
\gamma=\frac{\mathfrak E\sin2\Phi}{q^{\frac52}}
\]
\[
\begin{alignedat}{2}
{[}\mathfrak A&=\frac1{2aa(1-ee)\rho}
 &\qquad \log\mathfrak A&=1{,}407\,0739-10\\
\mathfrak B&=\frac{1-3ee}{6a^4(1-ee)^3\rho}
 &\log\mathfrak B&=7{,}317\,8248-30\\
\mathfrak C&=\frac{ee+15e^4}{1-3ee}
 &\log\mathfrak C&=7{,}868\,9880-10\\
\mathfrak D&=\frac{14e^4}{1-3ee}
 &\log\mathfrak D&=6{,}793\,4664-10\\
\mathfrak E&=\frac{ee}{2a^3(1-ee)\rho}
 &\log\mathfrak E&=2{,}424\,6792-20\\
\rho&=\frac1{206264{,}8..}\mathclose{]},
\end{alignedat}
\]
\noindent [wobei der Abplattungswert \(\dfrac1{302{,}68}\) verwendet wurde.
\(\alpha,\beta,\gamma\) gehören zum Argument
\(\tfrac13(x'+2x)=x+\tfrac13(x'-x)\)].

% Printed 160; principal PDF 167; Google PDF 171.
\par\clearpage
\noindent\makebox[0pt][l]{\fontsize{13}{16}\selectfont 160}\hfill{\small\textbf{NACHLASS.}}\hfill\null\par\medskip

\begin{center}
\(\left[\text{Tabelle für }\log\frac{1}{1-ee\sin\varphi^2}=\log q.\right]\)
\end{center}
{\footnotesize
\setlength{\tabcolsep}{4pt}
\renewcommand{\arraystretch}{1.03}
\noindent\makebox[\linewidth][c]{%
\begin{tabular}{|l|r|r|r||l|r|r|r|}
\hline
\multicolumn{1}{|c|}{\rule{0pt}{18pt}\shortstack{\(x\)\\(Meter)}} & \(\log\dfrac{10^7}{1-ee\sin\varphi^2}\) & \multicolumn{2}{c||}{Diff.}
 & \multicolumn{1}{c|}{\shortstack{\(x\)\\(Meter)}} & \(\log\dfrac{10^7}{1-ee\sin\varphi^2}\) & \multicolumn{2}{c|}{Diff.} \\[.5ex]
\cline{1-4}\cline{5-8}
\makebox[1.15em][r]{$+$}\makebox[1.65em][r]{800}\,000 & 14014,096 & \raisebox{-5pt}[0pt][0pt]{90,289} &  & \makebox[1.15em][r]{}\makebox[1.65em][r]{0} & 17597,264 & \raisebox{-5pt}[0pt][0pt]{87,869} & 129 \\
\makebox[1.15em][r]{}\makebox[1.65em][r]{78} & 104,385 & \raisebox{-5pt}[0pt][0pt]{90,298} & 9 & \makebox[1.15em][r]{$-$}\makebox[1.65em][r]{2} & 685,133 & \raisebox{-5pt}[0pt][0pt]{87,737} & 132 \\
\makebox[1.15em][r]{}\makebox[1.65em][r]{76} & 194,683 & \raisebox{-5pt}[0pt][0pt]{90,303} & 5 & \makebox[1.15em][r]{}\makebox[1.65em][r]{4} & 772,870 & \raisebox{-5pt}[0pt][0pt]{87,601} & 136 \\
\makebox[1.15em][r]{}\makebox[1.65em][r]{74} & 284,986 & \raisebox{-5pt}[0pt][0pt]{90,304} & 1 & \makebox[1.15em][r]{}\makebox[1.65em][r]{6} & 860,471 & \raisebox{-5pt}[0pt][0pt]{87,461} & 140 \\
\makebox[1.15em][r]{}\makebox[1.65em][r]{72} & 375,290 & \raisebox{-5pt}[0pt][0pt]{90,302} & 2 & \makebox[1.15em][r]{}\makebox[1.65em][r]{8} & 17947,932 & \raisebox{-5pt}[0pt][0pt]{87,319} & 142 \\
\makebox[1.15em][r]{$+$}\makebox[1.65em][r]{700}\,000 & 465,592 & \raisebox{-5pt}[0pt][0pt]{90,296} & 6 & \makebox[1.15em][r]{$-$}\makebox[1.65em][r]{100}\,000 & 18035,251 & \raisebox{-5pt}[0pt][0pt]{87,173} & 146 \\
\makebox[1.15em][r]{}\makebox[1.65em][r]{68} & 555,888 & \raisebox{-5pt}[0pt][0pt]{90,287} & 9 & \makebox[1.15em][r]{}\makebox[1.65em][r]{12} & 122,424 & \raisebox{-5pt}[0pt][0pt]{87,023} & 150 \\
\makebox[1.15em][r]{}\makebox[1.65em][r]{66} & 646,175 & \raisebox{-5pt}[0pt][0pt]{90,273} & 14 & \makebox[1.15em][r]{}\makebox[1.65em][r]{14} & 209,447 & \raisebox{-5pt}[0pt][0pt]{86,870} & 153 \\
\makebox[1.15em][r]{}\makebox[1.65em][r]{64} & 736,448 & \raisebox{-5pt}[0pt][0pt]{90,257} & 16 & \makebox[1.15em][r]{}\makebox[1.65em][r]{16} & 296,317 & \raisebox{-5pt}[0pt][0pt]{86,713} & 157 \\
\makebox[1.15em][r]{}\makebox[1.65em][r]{62} & 826,705 & \raisebox{-5pt}[0pt][0pt]{90,237} & 20 & \makebox[1.15em][r]{}\makebox[1.65em][r]{18} & 383,030 & \raisebox{-5pt}[0pt][0pt]{86,554} & 159 \\
\makebox[1.15em][r]{$+$}\makebox[1.65em][r]{600}\,000 & 14916,942 & \raisebox{-5pt}[0pt][0pt]{90,214} & 23 & \makebox[1.15em][r]{$-$}\makebox[1.65em][r]{200}\,000 & 469,584 & \raisebox{-5pt}[0pt][0pt]{86,391} & 163 \\
\makebox[1.15em][r]{}\makebox[1.65em][r]{58} & 15007,156 & \raisebox{-5pt}[0pt][0pt]{90,187} & 27 & \makebox[1.15em][r]{}\makebox[1.65em][r]{22} & 555,975 & \raisebox{-5pt}[0pt][0pt]{86,224} & 167 \\
\makebox[1.15em][r]{}\makebox[1.65em][r]{56} & 097,343 & \raisebox{-5pt}[0pt][0pt]{90,156} & 31 & \makebox[1.15em][r]{}\makebox[1.65em][r]{24} & 642,199 & \raisebox{-5pt}[0pt][0pt]{86,054} & 170 \\
\makebox[1.15em][r]{}\makebox[1.65em][r]{54} & 187,499 & \raisebox{-5pt}[0pt][0pt]{90,122} & 34 & \makebox[1.15em][r]{}\makebox[1.65em][r]{26} & 728,253 & \raisebox{-5pt}[0pt][0pt]{85,881} & 173 \\
\makebox[1.15em][r]{}\makebox[1.65em][r]{52} & 277,621 & \raisebox{-5pt}[0pt][0pt]{90,085} & 37 & \makebox[1.15em][r]{}\makebox[1.65em][r]{28} & 814,134 & \raisebox{-5pt}[0pt][0pt]{85,704} & 177 \\
\makebox[1.15em][r]{$+$}\makebox[1.65em][r]{500}\,000 & 367,706 & \raisebox{-5pt}[0pt][0pt]{90,043} & 42 & \makebox[1.15em][r]{$-$}\makebox[1.65em][r]{300}\,000 & 899,838 & \raisebox{-5pt}[0pt][0pt]{85,525} & 179 \\
\makebox[1.15em][r]{}\makebox[1.65em][r]{48} & 457,749 & \raisebox{-5pt}[0pt][0pt]{89,999} & 44 & \makebox[1.15em][r]{}\makebox[1.65em][r]{32} & 18985,363 & \raisebox{-5pt}[0pt][0pt]{85,341} & 184 \\
\makebox[1.15em][r]{}\makebox[1.65em][r]{46} & 547,748 & \raisebox{-5pt}[0pt][0pt]{89,950} & 49 & \makebox[1.15em][r]{}\makebox[1.65em][r]{34} & 19070,704 & \raisebox{-5pt}[0pt][0pt]{85,155} & 186 \\
\makebox[1.15em][r]{}\makebox[1.65em][r]{44} & 637,698 & \raisebox{-5pt}[0pt][0pt]{89,898} & 52 & \makebox[1.15em][r]{}\makebox[1.65em][r]{36} & 155,859 & \raisebox{-5pt}[0pt][0pt]{84,964} & 191 \\
\makebox[1.15em][r]{}\makebox[1.65em][r]{42} & 727,596 & \raisebox{-5pt}[0pt][0pt]{89,843} & 55 & \makebox[1.15em][r]{}\makebox[1.65em][r]{38} & 240,823 & \raisebox{-5pt}[0pt][0pt]{84,771} & 193 \\
\makebox[1.15em][r]{$+$}\makebox[1.65em][r]{400}\,000 & 817,439 & \raisebox{-5pt}[0pt][0pt]{89,784} & 59 & \makebox[1.15em][r]{$-$}\makebox[1.65em][r]{400}\,000 & 325,594 & \raisebox{-5pt}[0pt][0pt]{84,574} & 197 \\
\makebox[1.15em][r]{}\makebox[1.65em][r]{38} & 907,223 & \raisebox{-5pt}[0pt][0pt]{89,722} & 62 & \makebox[1.15em][r]{}\makebox[1.65em][r]{42} & 410,168 & \raisebox{-5pt}[0pt][0pt]{84,374} & 200 \\
\makebox[1.15em][r]{}\makebox[1.65em][r]{36} & 15996,945 & \raisebox{-5pt}[0pt][0pt]{89,656} & 66 & \makebox[1.15em][r]{}\makebox[1.65em][r]{44} & 494,542 & \raisebox{-5pt}[0pt][0pt]{84,171} & 203 \\
\makebox[1.15em][r]{}\makebox[1.65em][r]{34} & 16086,601 & \raisebox{-5pt}[0pt][0pt]{89,587} & 69 & \makebox[1.15em][r]{}\makebox[1.65em][r]{46} & 578,713 & \raisebox{-5pt}[0pt][0pt]{83,964} & 207 \\
\makebox[1.15em][r]{}\makebox[1.65em][r]{32} & 176,188 & \raisebox{-5pt}[0pt][0pt]{89,514} & 73 & \makebox[1.15em][r]{}\makebox[1.65em][r]{48} & 662,677 & \raisebox{-5pt}[0pt][0pt]{83,754} & 210 \\
\makebox[1.15em][r]{$+$}\makebox[1.65em][r]{300}\,000 & 265,702 & \raisebox{-5pt}[0pt][0pt]{89,437\rlap{.}} & 77 & \makebox[1.15em][r]{$-$}\makebox[1.65em][r]{500}\,000 & 746,431 & \raisebox{-5pt}[0pt][0pt]{83,540} & 214 \\
\makebox[1.15em][r]{}\makebox[1.65em][r]{28} & 355,139 & \raisebox{-5pt}[0pt][0pt]{89,357} & 80 & \makebox[1.15em][r]{}\makebox[1.65em][r]{52} & 829,971 & \raisebox{-5pt}[0pt][0pt]{83,324} & 216 \\
\makebox[1.15em][r]{}\makebox[1.65em][r]{26} & 444,496 & \raisebox{-5pt}[0pt][0pt]{89,274} & 83 & \makebox[1.15em][r]{}\makebox[1.65em][r]{54} & 913,295 & \raisebox{-5pt}[0pt][0pt]{83,104} & 220 \\
\makebox[1.15em][r]{}\makebox[1.65em][r]{24} & 533,770 & \raisebox{-5pt}[0pt][0pt]{89,187} & 87 & \makebox[1.15em][r]{}\makebox[1.65em][r]{56} & 19996,399 & \raisebox{-5pt}[0pt][0pt]{82,881} & 223 \\
\makebox[1.15em][r]{}\makebox[1.65em][r]{22} & 622,957 & \raisebox{-5pt}[0pt][0pt]{89,096} & 91 & \makebox[1.15em][r]{}\makebox[1.65em][r]{58} & 20079,280 & \raisebox{-5pt}[0pt][0pt]{82,655} & 226 \\
\makebox[1.15em][r]{$+$}\makebox[1.65em][r]{200}\,000 & 712,053 & \raisebox{-5pt}[0pt][0pt]{89,002} & 94 & \makebox[1.15em][r]{$-$}\makebox[1.65em][r]{600}\,000 & 161,935 & \raisebox{-5pt}[0pt][0pt]{82,425} & 230 \\
\makebox[1.15em][r]{}\makebox[1.65em][r]{18} & 801,055 & \raisebox{-5pt}[0pt][0pt]{88,905} & 97 & \makebox[1.15em][r]{}\makebox[1.65em][r]{62} & 244,360 & \raisebox{-5pt}[0pt][0pt]{82,193} & 232 \\
\makebox[1.15em][r]{}\makebox[1.65em][r]{16} & 889,960 & \raisebox{-5pt}[0pt][0pt]{88,803} & 102 & \makebox[1.15em][r]{}\makebox[1.65em][r]{64} & 326,553 & \raisebox{-5pt}[0pt][0pt]{81,957} & 236 \\
\makebox[1.15em][r]{}\makebox[1.65em][r]{14} & 16978,763 & \raisebox{-5pt}[0pt][0pt]{88,699} & 104 & \makebox[1.15em][r]{}\makebox[1.65em][r]{66} & 408,510 & \raisebox{-5pt}[0pt][0pt]{81,717} & 240 \\
\makebox[1.15em][r]{}\makebox[1.65em][r]{12} & 17067,462 & \raisebox{-5pt}[0pt][0pt]{88,591} & 108 & \makebox[1.15em][r]{}\makebox[1.65em][r]{68} & 490,227 & \raisebox{-5pt}[0pt][0pt]{81,474} & 243 \\
\makebox[1.15em][r]{$+$}\makebox[1.65em][r]{100}\,000 & 156,053 & \raisebox{-5pt}[0pt][0pt]{88,479} & 112 & \makebox[1.15em][r]{$-$}\makebox[1.65em][r]{700}\,000 & 571,701 & \raisebox{-5pt}[0pt][0pt]{81,230} & 244 \\
\makebox[1.15em][r]{}\makebox[1.65em][r]{8} & 244,532 & \raisebox{-5pt}[0pt][0pt]{88,364} & 115 & \makebox[1.15em][r]{}\makebox[1.65em][r]{72} & 652,931 & \raisebox{-5pt}[0pt][0pt]{80,980} & 250 \\
\makebox[1.15em][r]{}\makebox[1.65em][r]{6} & 332,896 & \raisebox{-5pt}[0pt][0pt]{88,246} & 118 & \makebox[1.15em][r]{}\makebox[1.65em][r]{74} & 733,911 & \raisebox{-5pt}[0pt][0pt]{80,728} & 252 \\
\makebox[1.15em][r]{}\makebox[1.65em][r]{4} & 421,142 & \raisebox{-5pt}[0pt][0pt]{88,124} & 122 & \makebox[1.15em][r]{}\makebox[1.65em][r]{76} & 814,639 & \raisebox{-5pt}[0pt][0pt]{80,473} & 255 \\
\makebox[1.15em][r]{}\makebox[1.65em][r]{2} & 509,266 & \raisebox{-5pt}[0pt][0pt]{87,998} & 126 & \makebox[1.15em][r]{}\makebox[1.65em][r]{78} & 895,112 & \raisebox{-5pt}[0pt][0pt]{80,214} & 259 \\
\makebox[1.15em][r]{}\makebox[1.65em][r]{0} & 17597,264 &  & 129 & \makebox[1.15em][r]{$-$}\makebox[1.65em][r]{800}\,000 & 20975,326 &  &  \\
\end{tabular}%
}\par}
\begin{center}
[Die Tabelle ist mit dem Abplattungswert \(\frac1{302{,}68}\) berechnet. Der
Ausgangspunkt für die \(x\)-Koordinaten hat die Breite
\(\varphi=51^\circ31^{\prime}48^{\prime\prime}{,}70\).]
\end{center}
% Printed 161; principal PDF 168; Google PDF 172.
\par\clearpage
\noindent\null\hfill{\small\textbf{KONFORME ABBILDUNG DES SPHÄROIDS IN DER EBENE.}}\hfill\makebox[0pt][r]{\fontsize{13}{16}\selectfont 161}\par\medskip

\begin{center}
[Tabelle zur Berechnung der Azimutkorrektur.]
\end{center}
{\scriptsize
\setlength{\tabcolsep}{1.6pt}
\renewcommand{\arraystretch}{1.08}
\noindent\makebox[\linewidth][c]{%
\begin{tabular}{|r|r|r|r|r||r|r|r|r|r|}
\hline
\multicolumn{1}{|c|}{\rule{0pt}{18pt}\shortstack{\(x\)\\(Meter)}} & \multicolumn{1}{c|}{\(\varphi\)} & \multicolumn{1}{c|}{\shortstack{\(\log\alpha\)\\\( 1{,}\ldots.-10\)}} & \multicolumn{1}{c|}{\shortstack{\(\log\beta\)\\\( 7{,}\ldots.-30\)}} & \multicolumn{1}{c||}{\shortstack{\(\log\gamma\)\\\( 2{,}\ldots.-20\)}} &
\multicolumn{1}{c|}{\shortstack{\(x\)\\(Meter)}} & \multicolumn{1}{c|}{\(\varphi\)} & \multicolumn{1}{c|}{\shortstack{\(\log\alpha\)\\\( 1{,}\ldots.-10\)}} & \multicolumn{1}{c|}{\shortstack{\(\log\beta\)\\\( 7{,}\ldots.-30\)}} & \multicolumn{1}{c|}{\shortstack{\(\log\gamma\)\\\( 2{,}\ldots.-20\)}} \\[.65ex]
\cline{1-5}\cline{6-10}
\rule{0pt}{1.5em}&&&&&&&&&\\
$+800\,000$ & \makebox[1.8em][r]{$44^\circ$}\makebox[1.5em][r]{$20^{\prime}$}\makebox[3.3em][r]{$1^{\prime\prime}{,}546$} & 1,404\,2711 & 7,31511 & 2,42106 & $+300\,000$ & \makebox[1.8em][r]{$48^\circ$}\makebox[1.5em][r]{$49^{\prime}$}\makebox[3.3em][r]{$57^{\prime\prime}{,}275$} & 1,403\,8208 & 7,31467 & 2,41671 \\
$790\,000$ & \makebox[1.8em][r]{}\makebox[1.5em][r]{$25$}\makebox[3.3em][r]{$25{,}585$} & 621 & 10 & 108 & $290\,000$ & \makebox[1.8em][r]{$48$}\makebox[1.5em][r]{$55$}\makebox[3.3em][r]{$21{,}061$} & 118 & 66 & 652 \\
$780\,000$ & \makebox[1.8em][r]{}\makebox[1.5em][r]{$30$}\makebox[3.3em][r]{$49{,}618$} & 530 & 09 & \makebox[0pt][r]{\(\mathord{{}^{\prime}_{.}}\)\,}109 & $280\,000$ & \makebox[1.8em][r]{$49$}\makebox[1.5em][r]{$0$}\makebox[3.3em][r]{$44{,}843$} & 1,403\,8029 & 65 & 633 \\
$770\,000$ & \makebox[1.8em][r]{}\makebox[1.5em][r]{$36$}\makebox[3.3em][r]{$13{,}646$} & 440 & 09 & 110 & $270\,000$ & \makebox[1.8em][r]{}\makebox[1.5em][r]{$6$}\makebox[3.3em][r]{$8{,}619$} & 1,403\,7939 & 64 & 611 \\
$760\,000$ & \makebox[1.8em][r]{}\makebox[1.5em][r]{$41$}\makebox[3.3em][r]{$37{,}669$} & 350 & 08 & 111 & $260\,000$ & \makebox[1.8em][r]{}\makebox[1.5em][r]{$11$}\makebox[3.3em][r]{$32{,}390$} & 850 & 63 & 590 \\
\cline{1-5}\cline{6-10}
$+750\,000$ & \makebox[1.8em][r]{}\makebox[1.5em][r]{$47$}\makebox[3.3em][r]{$1{,}687$} & 259 & 07 & 111 & $+250\,000$ & \makebox[1.8em][r]{}\makebox[1.5em][r]{$16$}\makebox[3.3em][r]{$56{,}157$} & 761 & 62 & 569 \\
$740\,000$ & \makebox[1.8em][r]{}\makebox[1.5em][r]{$52$}\makebox[3.3em][r]{$25{,}700$} & 169 & 06 & 110 & $240\,000$ & \makebox[1.8em][r]{}\makebox[1.5em][r]{$22$}\makebox[3.3em][r]{$19{,}918$} & 671 & 61 & 547 \\
$730\,000$ & \makebox[1.8em][r]{$44$}\makebox[1.5em][r]{$57$}\makebox[3.3em][r]{$49{,}708$} & 1,404\,2079 & 05 & 110 & $230\,000$ & \makebox[1.8em][r]{}\makebox[1.5em][r]{$27$}\makebox[3.3em][r]{$43{,}674$} & 582 & 60 & 525 \\
$720\,000$ & \makebox[1.8em][r]{$45$}\makebox[1.5em][r]{$3$}\makebox[3.3em][r]{$13{,}710$} & 1,404\,1988 & 04 & 108 & $220\,000$ & \makebox[1.8em][r]{}\makebox[1.5em][r]{$33$}\makebox[3.3em][r]{$7{,}426$} & 493 & 59 & 502 \\
$710\,000$ & \makebox[1.8em][r]{}\makebox[1.5em][r]{$8$}\makebox[3.3em][r]{$37{,}708$} & 898 & 03 & 107 & $210\,000$ & \makebox[1.8em][r]{}\makebox[1.5em][r]{$38$}\makebox[3.3em][r]{$31{,}172$} & 404 & 58 & 479 \\
\cline{1-5}\cline{6-10}
$+700\,000$ & \makebox[1.8em][r]{}\makebox[1.5em][r]{$14$}\makebox[3.3em][r]{$1{,}701$} & 808 & 02 & 105 & $+200\,000$ & \makebox[1.8em][r]{}\makebox[1.5em][r]{$43$}\makebox[3.3em][r]{$54{,}914$} & 315 & 58 & 455 \\
$690\,000$ & \makebox[1.8em][r]{}\makebox[1.5em][r]{$19$}\makebox[3.3em][r]{$25{,}689$} & 718 & 01 & 102 & $190\,000$ & \makebox[1.8em][r]{}\makebox[1.5em][r]{$49$}\makebox[3.3em][r]{$18{,}650$} & 226 & 57 & 431 \\
$680\,000$ & \makebox[1.8em][r]{}\makebox[1.5em][r]{$24$}\makebox[3.3em][r]{$49{,}671$} & 627 & 00 & 099 & $180\,000$ & \makebox[1.8em][r]{$49$}\makebox[1.5em][r]{$54$}\makebox[3.3em][r]{$42{,}381$} & 137 & 56 & 406 \\
$670\,000$ & \makebox[1.8em][r]{}\makebox[1.5em][r]{$30$}\makebox[3.3em][r]{$13{,}649$} & 537 & 7,31500 & 096 & $170\,000$ & \makebox[1.8em][r]{$50$}\makebox[1.5em][r]{$0$}\makebox[3.3em][r]{$6{,}108$} & 1,403\,7048 & 55 & 381 \\
$660\,000$ & \makebox[1.8em][r]{}\makebox[1.5em][r]{$35$}\makebox[3.3em][r]{$37{,}621$} & 447 & 7,31499 & 092 & $160\,000$ & \makebox[1.8em][r]{}\makebox[1.5em][r]{$5$}\makebox[3.3em][r]{$29{,}829$} & 1,403\,6959 & 54 & 356 \\
\cline{1-5}\cline{6-10}
$+650\,000$ & \makebox[1.8em][r]{}\makebox[1.5em][r]{$41$}\makebox[3.3em][r]{$1{,}589$} & 356 & 98 & 088 & $+150\,000$ & \makebox[1.8em][r]{}\makebox[1.5em][r]{$10$}\makebox[3.3em][r]{$53{,}546$} & 870 & 53 & 330 \\
$640\,000$ & \makebox[1.8em][r]{}\makebox[1.5em][r]{$46$}\makebox[3.3em][r]{$25{,}551$} & 266 & 97 & 084 & $140\,000$ & \makebox[1.8em][r]{}\makebox[1.5em][r]{$16$}\makebox[3.3em][r]{$17{,}258$} & 781 & 52 & 304 \\
$630\,000$ & \makebox[1.8em][r]{}\makebox[1.5em][r]{$51$}\makebox[3.3em][r]{$49{,}509$} & 176 & 96 & 079 & $130\,000$ & \makebox[1.8em][r]{}\makebox[1.5em][r]{$21$}\makebox[3.3em][r]{$40{,}964$} & 693 & 51 & 277 \\
$620\,000$ & \makebox[1.8em][r]{$45$}\makebox[1.5em][r]{$57$}\makebox[3.3em][r]{$13{,}461$} & 1,404\,1086 & 95 & 073 & $120\,000$ & \makebox[1.8em][r]{}\makebox[1.5em][r]{$27$}\makebox[3.3em][r]{$4{,}666$} & 604 & 51 & 250 \\
$610\,000$ & \makebox[1.8em][r]{$46$}\makebox[1.5em][r]{$2$}\makebox[3.3em][r]{$37{,}408$} & 1,404\,0995 & 94 & 067 & $110\,000$ & \makebox[1.8em][r]{}\makebox[1.5em][r]{$32$}\makebox[3.3em][r]{$28{,}362$} & 515 & 50 & 223 \\
\cline{1-5}\cline{6-10}
$+600\,000$ & \makebox[1.8em][r]{}\makebox[1.5em][r]{$8$}\makebox[3.3em][r]{$1{,}350$} & 905 & 93 & 061 & $+100\,000$ & \makebox[1.8em][r]{}\makebox[1.5em][r]{$37$}\makebox[3.3em][r]{$52{,}054$} & 427 & 49 & 195 \\
$590\,000$ & \makebox[1.8em][r]{}\makebox[1.5em][r]{$13$}\makebox[3.3em][r]{$25{,}288$} & 815 & 92 & 054 & $90\,000$ & \makebox[1.8em][r]{}\makebox[1.5em][r]{$43$}\makebox[3.3em][r]{$15{,}741$} & 338 & 48 & 166 \\
$580\,000$ & \makebox[1.8em][r]{}\makebox[1.5em][r]{$18$}\makebox[3.3em][r]{$49{,}220$} & 725 & 92 & 047 & $80\,000$ & \makebox[1.8em][r]{}\makebox[1.5em][r]{$48$}\makebox[3.3em][r]{$39{,}423$} & 250 & 47 & 137 \\
$570\,000$ & \makebox[1.8em][r]{}\makebox[1.5em][r]{$24$}\makebox[3.3em][r]{$13{,}147$} & 634 & 91 & 039 & $70\,000$ & \makebox[1.8em][r]{}\makebox[1.5em][r]{$54$}\makebox[3.3em][r]{$3{,}100$} & 162 & 46 & 108 \\
$560\,000$ & \makebox[1.8em][r]{}\makebox[1.5em][r]{$29$}\makebox[3.3em][r]{$37{,}069$} & 544 & 90 & 031 & $60\,000$ & \makebox[1.8em][r]{$50$}\makebox[1.5em][r]{$59$}\makebox[3.3em][r]{$26{,}772$} & 1,403\,6073 & 45 & 078 \\
\cline{1-5}\cline{6-10}
$+550\,000$ & \makebox[1.8em][r]{}\makebox[1.5em][r]{$35$}\makebox[3.3em][r]{$0{,}986$} & 454 & 89 & 023 & $+50\,000$ & \makebox[1.8em][r]{$51$}\makebox[1.5em][r]{$4$}\makebox[3.3em][r]{$50{,}439$} & 1,403\,5985 & 44 & 048 \\
$540\,000$ & \makebox[1.8em][r]{}\makebox[1.5em][r]{$40$}\makebox[3.3em][r]{$24{,}898$} & 364 & 88 & 014 & $40\,000$ & \makebox[1.8em][r]{}\makebox[1.5em][r]{$10$}\makebox[3.3em][r]{$14{,}101$} & 897 & 43 & 2,41017 \\
$530\,000$ & \makebox[1.8em][r]{}\makebox[1.5em][r]{$45$}\makebox[3.3em][r]{$48{,}805$} & 274 & 87 & 2,42005 & $30\,000$ & \makebox[1.8em][r]{}\makebox[1.5em][r]{$15$}\makebox[3.3em][r]{$37{,}758$} & 809 & 43 & 2,40986 \\
$520\,000$ & \makebox[1.8em][r]{}\makebox[1.5em][r]{$51$}\makebox[3.3em][r]{$12{,}707$} & 184 & 86 & 2,41995 & $20\,000$ & \makebox[1.8em][r]{}\makebox[1.5em][r]{$21$}\makebox[3.3em][r]{$1{,}410$} & 720 & 42 & 954 \\
$510\,000$ & \makebox[1.8em][r]{$46$}\makebox[1.5em][r]{$56$}\makebox[3.3em][r]{$36{,}603$} & 094 & 85 & 985 & $+10\,000$ & \makebox[1.8em][r]{}\makebox[1.5em][r]{$26$}\makebox[3.3em][r]{$25{,}058$} & 632 & 41 & 922 \\
\cline{1-5}\cline{6-10}
$+500\,000$ & \makebox[1.8em][r]{$47$}\makebox[1.5em][r]{$2$}\makebox[3.3em][r]{$0{,}495$} & 1,404\,0004 & 84 & 974 & $0$ & \makebox[1.8em][r]{}\makebox[1.5em][r]{$31$}\makebox[3.3em][r]{$48{,}700$} & 544 & 40 & 890 \\
$490\,000$ & \makebox[1.8em][r]{}\makebox[1.5em][r]{$7$}\makebox[3.3em][r]{$24{,}382$} & 1,403\,9914 & 83 & 963 & $-10\,000$ & \makebox[1.8em][r]{}\makebox[1.5em][r]{$37$}\makebox[3.3em][r]{$12{,}337$} & 457 & 39 & 857 \\
$480\,000$ & \makebox[1.8em][r]{}\makebox[1.5em][r]{$12$}\makebox[3.3em][r]{$48{,}264$} & 824 & 83 & 952 & $20\,000$ & \makebox[1.8em][r]{}\makebox[1.5em][r]{$42$}\makebox[3.3em][r]{$35{,}970$} & 369 & 38 & 823 \\
$470\,000$ & \makebox[1.8em][r]{}\makebox[1.5em][r]{$18$}\makebox[3.3em][r]{$12{,}140$} & 733 & 82 & 940 & $30\,000$ & \makebox[1.8em][r]{}\makebox[1.5em][r]{$47$}\makebox[3.3em][r]{$59{,}598$} & 281 & 37 & 790 \\
$460\,000$ & \makebox[1.8em][r]{}\makebox[1.5em][r]{$23$}\makebox[3.3em][r]{$36{,}012$} & 644 & 81 & 927 & $40\,000$ & \makebox[1.8em][r]{}\makebox[1.5em][r]{$53$}\makebox[3.3em][r]{$23{,}221$} & 193 & 36 & 755 \\
\cline{1-5}\cline{6-10}
$+450\,000$ & \makebox[1.8em][r]{}\makebox[1.5em][r]{$28$}\makebox[3.3em][r]{$59{,}879$} & 554 & 80 & 915 & $-50\,000$ & \makebox[1.8em][r]{$51$}\makebox[1.5em][r]{$58$}\makebox[3.3em][r]{$46{,}838$} & 106 & 36 & 721 \\
$440\,000$ & \makebox[1.8em][r]{}\makebox[1.5em][r]{$34$}\makebox[3.3em][r]{$23{,}740$} & 464 & 79 & 902 & $60\,000$ & \makebox[1.8em][r]{$52$}\makebox[1.5em][r]{$4$}\makebox[3.3em][r]{$10{,}451$} & 1,403\,5018 & 35 & 685 \\
$430\,000$ & \makebox[1.8em][r]{}\makebox[1.5em][r]{$39$}\makebox[3.3em][r]{$47{,}597$} & 374 & 78 & 888 & $70\,000$ & \makebox[1.8em][r]{}\makebox[1.5em][r]{$9$}\makebox[3.3em][r]{$34{,}059$} & 1,403\,4931 & 34 & 650 \\
$420\,000$ & \makebox[1.8em][r]{}\makebox[1.5em][r]{$45$}\makebox[3.3em][r]{$11{,}448$} & 284 & 77 & 874 & $80\,000$ & \makebox[1.8em][r]{}\makebox[1.5em][r]{$14$}\makebox[3.3em][r]{$57{,}663$} & 843 & 33 & 614 \\
$410\,000$ & \makebox[1.8em][r]{}\makebox[1.5em][r]{$50$}\makebox[3.3em][r]{$35{,}295$} & 194 & 76 & 859 & $90\,000$ & \makebox[1.8em][r]{}\makebox[1.5em][r]{$20$}\makebox[3.3em][r]{$21{,}261$} & 756 & 32 & 577 \\
\cline{1-5}\cline{6-10}
$+400\,000$ & \makebox[1.8em][r]{$47$}\makebox[1.5em][r]{$55$}\makebox[3.3em][r]{$59{,}136$} & 104 & 75 & 844 & $-100\,000$ & \makebox[1.8em][r]{}\makebox[1.5em][r]{$25$}\makebox[3.3em][r]{$44{,}855$} & 668 & 31 & 540 \\
$390\,000$ & \makebox[1.8em][r]{$48$}\makebox[1.5em][r]{$1$}\makebox[3.3em][r]{$22{,}973$} & 1,403\,9014 & 75 & 829 & $110\,000$ & \makebox[1.8em][r]{}\makebox[1.5em][r]{$31$}\makebox[3.3em][r]{$8{,}443$} & 581 & 30 & 503 \\
$380\,000$ & \makebox[1.8em][r]{}\makebox[1.5em][r]{$6$}\makebox[3.3em][r]{$46{,}804$} & 1,403\,8925 & 74 & 813 & $120\,000$ & \makebox[1.8em][r]{}\makebox[1.5em][r]{$36$}\makebox[3.3em][r]{$32{,}027$} & 494 & 29 & 465 \\
$370\,000$ & \makebox[1.8em][r]{}\makebox[1.5em][r]{$12$}\makebox[3.3em][r]{$10{,}631$} & 835 & 73 & 797 & $130\,000$ & \makebox[1.8em][r]{}\makebox[1.5em][r]{$41$}\makebox[3.3em][r]{$55{,}606$} & 407 & 29 & 426 \\
$360\,000$ & \makebox[1.8em][r]{}\makebox[1.5em][r]{$17$}\makebox[3.3em][r]{$34{,}452$} & 745 & 72 & 780 & $140\,000$ & \makebox[1.8em][r]{}\makebox[1.5em][r]{$47$}\makebox[3.3em][r]{$19{,}180$} & 320 & 28 & 387 \\
\cline{1-5}\cline{6-10}
$+350\,000$ & \makebox[1.8em][r]{}\makebox[1.5em][r]{$22$}\makebox[3.3em][r]{$58{,}268$} & 655 & 71 & 763 & $-150\,000$ & \makebox[1.8em][r]{}\makebox[1.5em][r]{$52$}\makebox[3.3em][r]{$42{,}749$} & 233 & 27 & 348 \\
$340\,000$ & \makebox[1.8em][r]{}\makebox[1.5em][r]{$28$}\makebox[3.3em][r]{$22{,}080$} & 566 & 70 & 746 & $160\,000$ & \makebox[1.8em][r]{$52$}\makebox[1.5em][r]{$58$}\makebox[3.3em][r]{$6{,}313$} & 146 & 26 & 308 \\
$330\,000$ & \makebox[1.8em][r]{}\makebox[1.5em][r]{$33$}\makebox[3.3em][r]{$45{,}886$} & 476 & 69 & 728 & $170\,000$ & \makebox[1.8em][r]{$53$}\makebox[1.5em][r]{$3$}\makebox[3.3em][r]{$29{,}872$} & 1,403\,4060 & 25 & 268 \\
$320\,000$ & \makebox[1.8em][r]{}\makebox[1.5em][r]{$39$}\makebox[3.3em][r]{$9{,}687$} & 387 & 68 & 710 & $180\,000$ & \makebox[1.8em][r]{}\makebox[1.5em][r]{$8$}\makebox[3.3em][r]{$53{,}427$} & 1,403\,3973 & 24 & 227 \\
$310\,000$ & \makebox[1.8em][r]{}\makebox[1.5em][r]{$44$}\makebox[3.3em][r]{$33{,}484$} & 297 & 67 & 693 & $190\,000$ & \makebox[1.8em][r]{}\makebox[1.5em][r]{$14$}\makebox[3.3em][r]{$16{,}977$} & 886 & 23 & 186 \\
\cline{1-5}\cline{6-10}
$+300\,000$ & \makebox[1.8em][r]{$48$}\makebox[1.5em][r]{$49$}\makebox[3.3em][r]{$57{,}275$} & 1,403\,8208 & 7,31467 & 2,41671 & $-200\,000$ & \makebox[1.8em][r]{$53$}\makebox[1.5em][r]{$19$}\makebox[3.3em][r]{$40{,}521$} & 1,403\,3800 & 7,31422 & 2,40145 \\
\end{tabular}%
}\par}
\medskip
\noindent\textsc{ix.}\hfill 21\hspace*{.12\linewidth}\par
\endgroup
